% SOURCE: Chinese PDF viewer pp. 141 (lower part)--144 (upper part).
\section{Symmetry and Conserved Quantities}
Define a transformation $\op S$ by
\[\ket\psi\longrightarrow\ket{\psi'}=\op S\ket\psi.\]
Before the transformation, $\ket\psi$ satisfies the Schr\"odinger equation
\[\ii\hbar\frac{\partial}{\partial t}\ket\psi=\op H\ket\psi.\]
Applying the time-independent transformation $\op S$ from the left gives
\[\ii\hbar\frac{\partial}{\partial t}\op S\ket\psi
=\op S\op H\ket\psi
=\op S\op H\op S^{-1}\op S\ket\psi.\]
After the transformation, $\ket{\psi'}$ satisfies the Schr\"odinger equation
\[\ii\hbar\frac{\partial}{\partial t}\ket{\psi'}
=\op H'\ket{\psi'},\qquad
\op H'=\op S\op H\op S^{-1}.\]
If $\op H$ remains unchanged after the transformation,
\[\op H'=\op S\op H\op S^{-1}=\op H,
\qquad\op S\op H=\op H\op S,\]
that is,
\[[\op S,\op H]=0,\]
then $\ket{\psi'}$ and $\ket\psi$ satisfy the same Schr\"odinger equation,
\[\ii\hbar\frac{\partial}{\partial t}\ket{\psi'}
=\op H\ket{\psi'}.\]
The system is said to be invariant, or symmetric, under $\op S$.

\textbf{Relationship between symmetry and conserved quantities.}

For a symmetry transformation $\op S$, if $\op S=\op S^{\dagger}$ and $\op S$ is a mechanical quantity, then $[\op S,\op H]=0$ shows that $\op S$ is a conserved mechanical quantity. If $\op S\ne\op S^{\dagger}$ but $\op S$ is generated by a mechanical quantity $\op F=\op F^{\dagger}$, i.e., $\op S=\op S(\op F)$, then $[\op S,\op H]=0$ implies $[\op F,\op H]=0$; thus $\op F$ is a conserved mechanical quantity.

\subsection*{1) Spatial-translation invariance and momentum conservation}
The spatial-translation operator is
\[\op S=e^{-\ii\op{\vect p}\cdot\vect x/\hbar},\]
generated by the momentum operator $\op{\vect p}$. If the system is spatially translation invariant, $[\op S,\op H]=0$, then $[\op{\vect p},\op H]=0$, leading to momentum conservation.

\subsection*{2) Time-translation invariance and energy conservation}
The time-translation operator is
\[\op S=e^{-\ii\op Ht/\hbar},\]
generated by the Hamiltonian $\op H$. If the system is time-translation invariant, $[\op S,\op H]=0$, then $[\op H,\op H]=0$, leading to energy conservation.

\subsection*{3) Spatial-rotation invariance and angular-momentum conservation}
The spatial-rotation operator is
\[\op S=e^{-\ii\op{\vect J}\cdot\vect e_n\,\varphi/\hbar},\]
generated by the angular-momentum operator $\op{\vect J}$. If the system is spatially rotation invariant, $[\op S,\op H]=0$, then $[\op{\vect J},\op H]=0$, leading to angular-momentum conservation.

\subsection*{4) Spatial-inversion invariance and parity conservation}
Define the spatial-inversion operator by
\[\vect r\longrightarrow\vect r'=\parityop\vect r\parityop^{-1}=-\vect r,
\qquad \psi(\vect r)\longrightarrow\psi'(\vect r)
=\parityop\psi(\vect r)=\psi(-\vect r).\]
It is a spatial inversion.

For arbitrary states $\psi(\vect r)$ and $\varphi(\vect r)$,
\begin{align*}
\int\dd^3r\,\psi^*(\vect r)(\parityop\varphi(\vect r))
&=\int\dd^3r\,\psi^*(\vect r)\varphi(-\vect r)\\
&=\int\dd^3r\,\psi^*(-\vect r)\varphi(\vect r)
=\int\dd^3r\,(\parityop\psi(\vect r))^*\varphi(\vect r).
\end{align*}
Hence $\parityop$ is a Hermitian operator and represents a mechanical quantity, called \emph{parity}.

Since
\[\parityop\psi(\vect r)=\psi(-\vect r),\qquad
\parityop^2\psi(\vect r)=\parityop\psi(-\vect r)=\psi(\vect r),\]
the eigenvalues of $\parityop^2$ are $1$. For the parity operator $\parityop$, the eigenvalues are $+1$ and $-1$, with eigenstates $\psi_s(\vect r)$ and $\psi_a(\vect r)$:
\[\parityop\psi_s(\vect r)=\psi_s(\vect r),\quad
\parityop\psi_a(\vect r)=-\psi_a(\vect r).\]
They are respectively even and odd states,
\[\psi_s(\vect r)=\psi_s(-\vect r),\qquad
\psi_a(\vect r)=-\psi_a(-\vect r).\]
If the system possesses spatial-inversion invariance, $[\parityop,\op H]=0$, then
\begin{enumerate}[label=\alph*)]
\item parity is conserved;
\item $\parityop$ and $\op H$ have common eigenstates.
\end{enumerate}
\textbf{Example.} For $\op H=-\hbar^2\nabla^2/(2m)+V(\vect r)$, if $V(-\vect r)=V(\vect r)$, then $\parityop\op H\parityop^{-1}=\op H$; hence parity is conserved and the Hamiltonian has common eigenstates with $\parityop$. It can be shown that, for a one-dimensional bound-state problem with no degeneracy, all eigenstates of $\op H$ are also eigenstates of $\parityop$.

\textbf{Example.} Let
\[\op{\vect L}=\vect r\times\op{\vect p},\qquad
\op{\vect p}=-\ii\hbar\vect\nabla.\]
Then
\[\parityop\op{\vect L}\parityop^{-1}
=(\parityop\vect r\parityop^{-1})\times(\parityop\op{\vect p}\parityop^{-1})
=(-\vect r)\times(-\op{\vect p})=\op{\vect L},\]
and
\[[\parityop,\op{\vect L}]=0,\qquad
[\parityop,\op{\vect L}^{\,2}]=0.\]
Thus $\parityop$, $\op{\vect L}^{\,2}$ and $\op L_z$ have common eigenstates. The common eigenstates of $\op{\vect L}^{\,2}$ and $\op L_z$ are the spherical harmonics $Y_{lm}(\theta,\varphi)$:
\[Y_{lm}(\theta,\varphi)=N_{lm}P_l^{|m|}(\cos\theta)e^{\ii m\varphi},\]
where
\[P_l^{|m|}(\cos\theta)=\frac1{2^l l!}
(1-\cos^2\theta)^{|m|/2}
\left(\frac{\dd}{\dd\cos\theta}\right)^{l+|m|}
(\cos^2\theta-1)^l.\]
Under the inversion $\parityop$, $\vect r\to-\vect r$,
\[r\to r,\qquad\theta\to\pi-\theta,\qquad
\varphi\to\pi+\varphi,\qquad
\cos\theta\to-\cos\theta.\]
Therefore,
\[P_l^{|m|}(\cos\theta)\to(-1)^{l+|m|}P_l^{|m|}(\cos\theta),
\qquad e^{\ii m\varphi}\to(-1)^m e^{\ii m\varphi},\]
so that
\[\parityop\,Y_{lm}(\theta,\varphi)
=(-1)^{l+2|m|}Y_{lm}(\theta,\varphi)=(-1)^lY_{lm}(\theta,\varphi).\]
The common eigenstates of $\parityop$, $\op{\vect L}^{\,2}$ and $\op L_z$ are the spherical harmonics $Y_{lm}(\theta,\varphi)$, and the eigenvalue of $\parityop$ is $(-1)^l$.

\textbf{Question.} Why is there no physical quantity called parity in classical mechanics?

\textbf{Answer.} In classical mechanics there is no interference, and one cannot distinguish $\vect r$ from $-\vect r$.